\documentclass[10pt,a4paper]{amsart}
\usepackage{amsmath,amssymb,amsthm,graphicx}
\usepackage{verbatim}
\newenvironment{reference}{\comment}{\endcomment}
%\newenvironment{reference}{}{}
\newenvironment{slogan}{\comment}{\endcomment}
\newenvironment{history}{\comment}{\endcomment}
% For commutative diagrams we use Xy-pic
\usepackage[all]{xy}

% We use 2cell for 2-commutative diagrams.
\xyoption{2cell}
% Theorem environments.
%
\theoremstyle{plain}
\newtheorem{theorem}[subsection]{Theorem}
\newtheorem{proposition}[subsection]{Proposition}
\newtheorem{lemma}[subsection]{Lemma}

\theoremstyle{definition}
\newtheorem{definition}[subsection]{Definition}
\newtheorem{example}[subsection]{Example}
\newtheorem{exercise}[subsection]{Exercise}
\newtheorem{situation}[subsection]{Situation}

\theoremstyle{remark}
\newtheorem{remark}[subsection]{Remark}
\newtheorem{remarks}[subsection]{Remarks}

\numberwithin{equation}{subsection}

% Macros
%
\def\lim{\mathop{\mathrm{lim}}\nolimits}
\def\colim{\mathop{\mathrm{colim}}\nolimits}
\def\Spec{\mathop{\mathrm{Spec}}}
\def\Hom{\mathop{\mathrm{Hom}}\nolimits}
\def\Ext{\mathop{\mathrm{Ext}}\nolimits}
\def\SheafHom{\mathop{\mathcal{H}\!\mathit{om}}\nolimits}
\def\SheafExt{\mathop{\mathcal{E}\!\mathit{xt}}\nolimits}
\def\Sch{\mathit{Sch}}
\def\Mor{\mathop{\mathrm{Mor}}\nolimits}
\def\Ob{\mathop{\mathrm{Ob}}\nolimits}
\def\Sh{\mathop{\mathit{Sh}}\nolimits}
\def\NL{\mathop{N\!L}\nolimits}
\def\CH{\mathop{\mathrm{CH}}\nolimits}
\def\proetale{{pro\text{-}\acute{e}tale}}
\def\etale{{\acute{e}tale}}
\def\QCoh{\mathit{QCoh}}
\def\Ker{\mathop{\mathrm{Ker}}}
\def\Im{\mathop{\mathrm{Im}}}
\def\Coker{\mathop{\mathrm{Coker}}}
\def\Coim{\mathop{\mathrm{Coim}}}

% Boxtimes
%
\DeclareMathSymbol{\boxtimes}{\mathbin}{AMSa}{"02}

%
% Macros for moduli stacks/spaces
%
\def\QCohstack{\mathcal{QC}\!\mathit{oh}}
\def\Cohstack{\mathcal{C}\!\mathit{oh}}
\def\Spacesstack{\mathcal{S}\!\mathit{paces}}
\def\Quotfunctor{\mathrm{Quot}}
\def\Hilbfunctor{\mathrm{Hilb}}
\def\Curvesstack{\mathcal{C}\!\mathit{urves}}
\def\Polarizedstack{\mathcal{P}\!\mathit{olarized}}
\def\Complexesstack{\mathcal{C}\!\mathit{omplexes}}
% \Pic is the operator that assigns to X its picard group, usage \Pic(X)
% \Picardstack_{X/B} denotes the Picard stack of X over B
% \Picardfunctor_{X/B} denotes the Picard functor of X over B
\def\Pic{\mathop{\mathrm{Pic}}\nolimits}
\def\Picardstack{\mathcal{P}\!\mathit{ic}}
\def\Picardfunctor{\mathrm{Pic}}
\def\Deformationcategory{\mathcal{D}\!\mathit{ef}}
\title{Quasi-compactness of a morphism of algebraic stacks can be checked after a flat cover of the target}
\author{The Stacks Project Authors}
\date{stacks-project:stacks-morphisms.tex:local-snapshot-20261012; GNU FDL 1.2 or later}
\begin{document}
\maketitle
\noindent\textit{Editorial note.} All delivered mathematical text is the author's own, copied from the pinned Stacks Project chapter source sources/stacks-algebraic-stacks-morphisms/pinned-original-source.tex; no statement or proof step was rewritten, reordered or removed.\par\medskip
\noindent\textit{Editorial note.} Preamble provenance: the Stacks Project keeps one single shared preamble for every chapter; this card takes it directly from the staged sources/stacks-algebraic-stacks-morphisms/preamble.tex as three verbatim spans (lines 8-12, 14-18 and 154-217), so no derived file was written and nothing was added or redefined. The wrapper class and the base packages are supplied by the assembler.\par\medskip
\noindent\textit{Editorial note.} Bibliography provenance: neither the counted statement nor its proof contains a citation (this lane rescans the whole span before writing the card), so the card carries no bibliography and no entry was reconstructed.\par\medskip
\noindent\textit{Editorial note.} Selection provenance: this card was opened under the file-based control contract (CONTROL.md), with the mechanical pre-card gate run on its statement and proof interval before the card was written.\par\medskip
\noindent\textit{Allowed standard prerequisites (editorial).}\begin{itemize}
\item The definition of a quasi-compact morphism of algebraic stacks through the quasi-compactness of the base changes along morphisms from quasi-compact algebraic stacks. The author's section on quasi-compact morphisms, in which the counted statement occurs. Formulates what has to be proved after the given base change is known to be quasi-compact.
\item The existence of a smooth surjection from a scheme onto an algebraic stack, and the possibility of shrinking a flat surjective cover of a quasi-compact stack to a quasi-compact open. Used inside the counted proof (source lines 2165-2171). Produces a quasi-compact scheme cover of the test stack through which the quasi-compactness can be transported.
\end{itemize}\medskip
\begin{lemma}
\label{lemma-check-quasi-compact-covering}
Let $f : \mathcal{X} \to \mathcal{Y}$ be a morphism of algebraic stacks.
Let $W \to \mathcal{Y}$ be surjective, flat, and locally of finite
presentation where $W$ is an algebraic space. If the base change
$W \times_\mathcal{Y} \mathcal{X} \to W$ is quasi-compact, then
$f$ is quasi-compact.
\end{lemma}

\begin{proof}
Assume $W \times_\mathcal{Y} \mathcal{X} \to W$ is quasi-compact.
Let $\mathcal{Z} \to \mathcal{Y}$ be a morphism with $\mathcal{Z}$
a quasi-compact algebraic stack. Choose a scheme $U$ and a surjective
smooth morphism
$U \to W \times_\mathcal{Y} \mathcal{Z}$.
Since $U \to \mathcal{Z}$ is flat, surjective, and locally of finite
presentation and $\mathcal{Z}$ is quasi-compact, we can find a
quasi-compact open subscheme $U' \subset U$ such that
$U' \to \mathcal{Z}$ is surjective.
Then
$U' \times_\mathcal{Y} \mathcal{X} =
U' \times_W (W \times_\mathcal{Y} \mathcal{X})$
is quasi-compact by assumption and surjects onto
$\mathcal{Z} \times_\mathcal{Y} \mathcal{X}$.
Hence $\mathcal{Z} \times_\mathcal{Y} \mathcal{X}$
is quasi-compact as desired.
\end{proof}

\end{document}
